\section{Generalization}

In this section, we generalize the results for the Fibonacci polynomials to $W_n(T)$ \emph{with $\deg f(T) = 1$ and $\deg g(T) = 0$}, including the Chebyshev polynomials of the second kind (Example \ref{ex:chebyshev2}). We also find all the polynomials in the Lucas-type sequence $w_n(T)$ which are perfect powers. 

First, we need the following factorization lemmas, similar to Lemma \ref{lem:Fpk} and \ref{lem:Fpkm}.
We skip the details since the proofs are similar.

\begin{lemma}
\label{lem:Wwpk}
    Let $p$ be a prime and $n = p^k$ for some $k \in \mathbb{N}$. Then $W_n(T) = \Delta(T)^{n-1}$ and $w_n(T) = f(T)^n$ in $\bF_q[T]$.
\end{lemma}


\begin{lemma}
\label{lem:Wwpkm}
    Let $n = p^km$ for some $k \in \mathbb{N}$. Then, $W_n(T) = W_{p^k}(T) W_{m} (T)^{p^k} = \Delta(T)^{p^k - 1} W_{m}(T)^{p^k}$ and $w_n(T) = w_m(T)^{p^k}$ in $\mathbb{F}_q[T]$.
    In particular, $w_n(T)$ is a perfect power if and only if it is powerful.
\end{lemma}


We also have the following coprimality result analogous to Corollary \ref{cor:coprime}.
\begin{proposition}
\label{prop:Wmcoprime}
    Suppose $p$ is an odd prime and $p \nmid m \cdot \lc(f) \cdot g$. Then $f(T)^2 + 4g(T)$ and $W_m(T)$ are coprime in $\mathbb{F}_q[T]$.
\end{proposition}

\begin{proof}
    Suppose $f(T)^2 + 4g(T)$ and $W_m(T)$ in $\mathbb{F}_q[T]$ share a factor. Then, there must exist some common zero of the polynomial, which we will denote as $\alpha \in \bF_{q^2}$.
    Then $f(\alpha)^2 + 4g(\alpha) = 0$ and \(W_m(\alpha) = 0\).
    However, such \(\alpha\) cannot exist by Lemma \ref{lem:florez_mod}.
\end{proof}

We will now let $\deg f = 1$ and $\deg g = 0$ in the recursive definition of $W_n(T)$ and $w_n(T)$ (see Definition \ref{lucassq}). Using the discriminant formulas in \cite{florez2019resultant}, we can show the following results analogous to the case of Fibonacci polynomials.

\begin{proposition}
    \label{prop:Wwsqfree}
    Suppose $p$ is odd and $p \nmid m \cdot \lc(f) \cdot g$, where $\lc(f)$ is the leading coefficient of $f(T)$. Then, $W_m(T)$ and $w_m(T)$ are square-free over $\mathbb{F}_q$.
\end{proposition}

\begin{proof}
    This follows from the discriminant formulas for $W_m$ and $w_m$ \cite[Theorem 4, Theorem 5]{florez2019resultant}:
    \begin{align*}
        \mathrm{Disc}(W_m) &= (-g)^{\frac{(m-1)(m-2)}{2}} 2^{m-1} m^{m-3} \lc(f)^{(m-1)(m-2)}, \\
        \mathrm{Disc}(w_m) &= (-g)^{\frac{m(m-1)}{2}} 2^{m-1} m^{m} \lc(f)^{m(m-1)} \text{.}
    \end{align*}
    In particular, the discriminants are nonzero for $p \nmid m \cdot \lc(f) \cdot g$, and $W_m$ and $w_m$ are thus square-free over $\mathbb{F}_q$.
\end{proof}



\begin{theorem}
    Let $p$ be an odd prime number and assume $p \nmid \lc(f) \cdot g$. If $j > 1$ and $p \mid j$, then there exists no $n$ such that $W_n(T)$ is a perfect $j$-th power in $\mathbb{F}_q[T]$.
\end{theorem}

\begin{proof}
    One can proceed as in the proof of Theorem \ref{thm:Fnnonpower}.
    Suppose $W_{p^km} (T)$ is a perfect $j$-th power for some $k \ge 0$ and $m$ not divisible by $p$.
    If $k = 0$, $W_{p^k m}(T) = W_{m}(T)$ is square-free by Proposition \ref{prop:Wwsqfree}, which cannot be a perfect $j$-th power for $j > 1$.
    Hence $k \ge 1$, and Lemma \ref{lem:Wwpkm} gives $W_{p^k m}(T) = (f(T)^2 + 4g(T))^{\frac{p^k - 1}{2}} W_m(T)^{p^k}$, where the coprimality of $f^2 + 4g$ and $W_m$ (Proposition \ref{prop:Wmcoprime}) implies that the exponent $\frac{p^k - 1}{2}$ has to be a multiple of $j$.
    This contradicts to $p \mid j$.
\end{proof}

\begin{theorem}
    \label{thm:Wn_perfect_power}
    Let $j > 1$. Let $p$ be an odd prime number coprime to $2j$, and let $a = \mathrm{ord}_{2j}(p)$ denote the order of $p$ in $\left( \mathbb{Z}/2j\mathbb{Z} \right)^\times$. Suppose $p \nmid \lc(f) \cdot g$. For $n > 0$, $W_n(T)$ is a perfect $j$-th power in $\mathbb{F}_q[T]$ if and only if $n = p^{ak}$ for some $k \in \mathbb{N}$.
\end{theorem}

\begin{proof}
    $(\Leftarrow)$ Let $k \in \mathbb{N}$ be arbitrary. By Lemma \ref{lem:Wwpk}, $W_{p^{ak}}(T) = \Delta(T)^{p^{ak} - 1} = \left(f(T)^2 + 4 g(T) \right)^{\frac{p^{ak} - 1}{2}}$. Because $p^a \equiv 1 \pmod{2j}$ and thus $\frac{p^{ak} - 1}{2} \equiv 0 \mod j$, $W_{p^{ak}}(T)$ is a perfect $j$-th power.
    
    $(\Rightarrow)$ Suppose $W_n(T)$ is a perfect $j$-th power in $\mathbb{F}_q[T]$ and write $n = p^{ak + b}m$, where $p$ and $m$ are coprime (so $p \nmid m$), $k \geq 0$, and \(0 \le b < a\). With Lemma \ref{lem:Wwpkm}, we may write $W_n(T)$ as follows:
    \begin{align*}
        W_{p^{ak + b}m}(T) &= W_{p^{ak + b}}(T) \left(W_m(T)\right)^{p^{ak + b}} = \left(f(T)^2 + 4 g(T) \right)^{\frac{p^{ak +b} - 1}{2}} \left(W_m(T)\right)^{p^{ak + b}}
    \end{align*}

    Suppose $b > 0$. Then, $p^{ak + b} \not \equiv 1 \pmod{2j}$ by the assumption $b < a$ and thus $\frac{p^{ak + b} - 1}{2} \not \equiv 0 \pmod j$. So $W_{p^{ak + b}}(T)$ is not a perfect $j$-th power. Then, $W_{p^{ak + b}m}(T)$ is not a perfect $j$-th power since $W_{p^{ak + b}}(T)$ and $W_m(T)$ are coprime by Proposition \ref{prop:Wmcoprime}. We thus let $b = 0$ and rewrite $W_n(T)$ as follows: $$W_{p^{ak}m}(T) = W_{p^{ak}}(T) \left(W_m(T)\right)^{p^{ak}} \text{,}$$ where $W_{p^{ak}}(T)$ is a perfect $j$-th power and $W_m(T)$ is square-free by Proposition \ref{prop:Wwsqfree}. Then, if $m \neq 1$, then $W_m(T) \neq 1$ and by Proposition \ref{prop:Wmcoprime} $j > 1$ must divide $p^{ak}$, which is a contradiction to $\gcd(p,2j) = 1$. We thus have $m = 1$ and $n = p^{ak}$.
\end{proof}


\begin{theorem}
    \label{thm:Wn_powerful}
    Suppose $p > 3$ and assume $p \nmid \lc(f) \cdot g$.
    Then \(W_n(T)\) is powerful in \(\bF_q[T]\) if and only if $n$ is a multiple of $p$.
\end{theorem}

\begin{proof}
    $(\Rightarrow)$ If $p \nmid n$, then $W_n(T)$ is square-free and cannot be a powerful polynomial.

    $(\Leftarrow)$ If $n = p^k m$ with $k \ge 1$, then $W_n(T)$ is powerful by Lemma \ref{lem:Wwpkm}.
\end{proof}

When $p = 2$, we have the following results, generalizing Theorem \ref{thm:mainFnp2} and Corollary \ref{cor:powerfulp2}.
\begin{lemma}
    \label{lem:Wnexpansion}
    For each $n$, there exist integers $\{c_{n, k}\}_{0 \le k \le \frac{n-1}{2}}$ such that
    \begin{equation}
        \label{eqn:Wnexpansion}
        W_n = \sum_{k=0}^{\lfloor\frac{n-1}{2}\rfloor} c_{n, k} f^{n-1 - 2k} g^k
    \end{equation}
    where $c_{n, (n-1)/2} = 1$ for odd $n$.
    In particular, $f \mid W_{n}$ for even $n$.
    When $n$ is odd and $g$ is a nonzero constant, then $\gcd(W_{n}, f) = 1$.
\end{lemma}
\begin{proof}
    The equation \eqref{eqn:Wnexpansion} easily follows from induction on $n$, and the following claims immediately follow from it.
\end{proof}

\begin{lemma}
    \label{lem:Wnsquare}
    Let $q$ be a power of 2 and $g \in \bZ$. Over $\bF_q$, we have
    \begin{enumerate}
        \item $W_{2n}(T) = f(T) W_n(T)^2$
        \item $W_{2n+1}(T) = (W_{n+1}(T) + g W_{n}(T))^2$.
    \end{enumerate}
\end{lemma}
\begin{proof}
    The proof for $F_n(T)$ in \cite[Lemma 3.5]{kitayama2017irreducibility} generalizes.
    The first claim follows from the Binet form \eqref{eqn:Wnbinet} and $\Delta = f$. 
    For the second claim, we have $f(T) W_{2n+1}(T) = W_{2n+2}(T) - g W_{2n}(T) = f(T)(W_{n+1}(T) + g W_{n}(T))^2$ (here we used $g^2 \equiv g \equiv -g\pmod{2}$), and the result follows.
\end{proof}

\begin{theorem}
    Let $p = 2$ and $q$ be a power of 2.
    Let $j > 2$ be an odd number and $a = \mathrm{ord}_j(2)$ be the order of 2 in $(\bZ / j\bZ)^\times$.
    Then $W_n(T)$ is a perfect $j$-th power in $\bF_q[T]$ if and only if $n = 2^{ak}$ for some $k \ge 1$.
\end{theorem}
\begin{proof}
    It can be proved similarly as Theorem \ref{thm:mainFnp2}.
    For $n = 2^k m$ with odd $m$, note that $W_n(T) = f(T)^{2^k - 1} W_m(T)^{2^k}$ by Lemma \ref{lem:Wwpkm} and \ref{lem:Wnexpansion}.
\end{proof}

\begin{theorem}
    Let $p = 2$ and $q$ be a power of 2.
    $W_n(T)$ is powerful over $\bF_q$ if and only if $n \ge 3$ is odd or a multiple of $4$.
\end{theorem}
\begin{proof}
    The proof is similar to Corollary \ref{cor:powerfulp2}, where we use Lemma \ref{lem:Wwpkm} and \ref{lem:Wnsquare} instead.
\end{proof}

\begin{theorem}
    Let $p = 3$ and $q$ be a power of 3.
    $W_n(T)$ is powerful over $\bF_q$ if and only if $9 \mid n$.
\end{theorem}
\begin{proof}
    The proof is similar to Corollary \ref{cor:powerfulp3}.
\end{proof}

\begin{example}
\label{ex:chebyshev2}
    When $f(T) = 2T$ and $g(T) = -1$, we have $W_n(T) = U_{n-1}(T)$, the Chebyshev polynomial of the second kind.
    By Theorem \ref{thm:Wn_perfect_power} and \ref{thm:Wn_powerful}, when $p$ is odd, $U_n(T)$ is a perfect power (resp. powerful) in $\bF_q[T]$ if and only if $n = p^{\mathrm{ord}_{2j}(p)k} - 1$ for some $k \in \bN$ (resp. $n \equiv -1 \pmod{p}$ for $p > 3$).
    Table \ref{tab:chebyshev2} shows factorizations of $U_n(T)$ over $\bF_3$ and $\bF_5$ for $1 \le n \le 20$.
    Note that $U_{2m}(T) \equiv 1 \pmod{2}$ and $U_{2m+1}(T) \equiv 0 \pmod{2}$ for all $m \ge 0$.
\end{example}


Similarly, we can characterize when $w_n(T)$ is a perfect power.
\begin{corollary}
\label{cor:wn_perfect_power}
    Let $p$ be an odd prime and assume $p \nmid \lc (f) \cdot g$. Then for $n > 0$, $w_n(T)$ is a perfect power in $\mathbb{F}_q[T]$ if and only if $n$ is a multiple of $p$.
\end{corollary}
\begin{proof}
    If $n = p^k m$ for some positive integer $k$ and $m$, $w_n(T)$ is a perfect power by Lemma \ref{lem:Wwpkm}.
    
    Suppose $n$ is not divisible by $p$. By Proposition \ref{prop:Wwsqfree}, $w_n(T)$ is square-free, and thus $w_n(T)$ is not a perfect power.
\end{proof}

\begin{example}
    When $f(T) = T$ and $g(T) = 1$, $w_n(T)$ recovers the Lucas polynomial $L_n(T)$.
    Thus $L_n(T)$ is a perfect power in $\bF_q[T]$ if and only if $p \mid n$.
    Table \ref{tab:lucas} shows factorizations of $L_n(T)$ over $\bF_2$, $\bF_3$, and $\bF_5$ for $1 \le n \le 20$.
\end{example}

\begin{example}
    When $f(T) = 2T$ and $g(T) = -1$, we have $w_n(T) = 2 T_n(T)$ where $T_n(T)$ is the Chebyshev polynomial of the first kind.
    Thus $T_n(T)$ is a perfect power in $\bF_q[T]$ if and only if $p \mid n$.
    Table \ref{tab:chebyshev1} shows factorizations of $T_n(T)$ over $\bF_3$ and $\bF_5$ for $1 \le n \le 20$.
    Note that Bhargava and Zieve studied factorizations of Dickson polynomials $D_n(T, a)$ over finite fields as well \cite{bhargava1999factoring}, which recovers the Chebyshev polynomial as $D_n(2T, 1) = 2 T_n(T)$.
\end{example}


When $f(T) = 1$ and $g(T) = 2T$, we have $W_n(T) = J_n(T)$, the Jacobsthal polynomial \cite{horadam1997jacobsthal}.
Unfortunately, Theorem \ref{thm:Wn_perfect_power} and Theorem \ref{thm:Wn_powerful} do not apply to these polynomials, since Florez--Higuita--Ramirez's discriminant formula \cite[Theorem 4]{florez2019resultant} does not cover the case of $\deg f = 0$ and $\deg g = 1$.
However, we can still determine when $J_n(T)$ is a perfect power or powerful in $\bF_q[T]$, by using the idea of the second proof of Proposition \ref{prop:Fmsqfree}.
Since $J_n(T) \equiv 1 \pmod{2}$ for all $n$, we will assume $p$ is odd.

From the Binet's formula \eqref{eqn:Wnbinet}, we have
\[
J_n(T) = \frac{\left(\frac{1 + \sqrt{1 + 8T}}{2}\right)^n - \left(\frac{1 - \sqrt{1 + 8T}}{2}\right)^n}{\sqrt{1 + 8T}}, \quad \alpha(T) = \frac{1 + \sqrt{1 + 8T}}{2}, \quad \beta(T) = \frac{1 - \sqrt{1 + 8T}}{2}.
\]
We have $\deg J_n = \lceil \frac{n}{2} \rceil - 1$, and the above formula tells us that the zeros $x$ of $J_n$ satisfies $\alpha(x)^n = \beta(x)^n \Leftrightarrow (\alpha(x) / \beta(x))^n = 1$ with $\alpha(x) \ne \beta(x)$.
Thus $x$ has a form of
\[
\frac{1 + \sqrt{1 + 8x}}{1 - \sqrt{1 + 8x}} = \zeta_n^k \Leftrightarrow x = x_k = - \frac{\zeta_n^k}{2(1 + \zeta_n^k)^2},\quad k = 1, 2, \dots, \left \lceil \frac{n}{2} \right \rceil - 1.
\]
One can check that $2x_k$ are algebraic integers for all $k$.
In particular, we can reduce $x_k$ modulo $p$ and get elements $\overline{x}_k \in \overline{\bF}_p$.
Using Lemma \ref{distinct} and $x / (1 + x)^2 = y / (1 + y)^2 \Leftrightarrow (x - y)(1 - xy) = 0$, we can show that all the zeros $\overline{x}_k$ for $k = 1, 2, \dots, \lceil \frac{n}{2} \rceil - 1$ are distinct in $\overline{\bF}_p$ when $p \nmid n$. Hence,
\begin{proposition}
    \label{thm:Jmsqfree}
    Let $p$ be an odd prime and $p \nmid m$. Then $J_m(T)$ is square-free in $\bF_q[T]$.
\end{proposition}

Now, using Lemma \ref{lem:Wwpk}, Lemma \ref{lem:Wwpkm}, and $J_1 = J_2 = 1$, we can characterize the perfect powers and powerful polynomials among $J_n(T)$ over $\bF_q$.
Since the argument is similar to that of Theorem \ref{thm:Wn_perfect_power} and \ref{thm:Wn_powerful}, we omit the details.
\begin{theorem}
\label{thm:Jn}
    Let $p$ be an odd prime and $q$ be a power of $p$.
    Let $j > 1$ and $a = \mathrm{ord}_{2j}(p)$ be the order of $p$ in $(\bZ / 2j \bZ)^\times$.
    Then for $n > 2$, $J_n(T)$ is a perfect $j$-th power in $\bF_q[T]$ if and only if $n = p^{ak}$ or $n = 2p^{ak}$ for some $k \ge 1$.
    Also, $J_n(T)$ is powerful in $\bF_q[T]$ if and only if $p \mid n$ when $p > 3$, or $9 \mid n$ when $p = 3$.
\end{theorem}
Table \ref{tab:jacobsthal} shows factorizations of $J_n(T)$ over $\bF_3$ and $\bF_5$ for $1 \le n \le 27$.
It would be interesting to generalize Theorem \ref{thm:Jn} to other Fibonacci-type polynomials $W_n(T)$ with $\deg f = 0$ and $\deg g = 1$.



\begin{table}
\centering
\begin{tabular}{c|p{0.30\textwidth}|p{0.30\textwidth}|p{0.30\textwidth}}
\toprule
$n$ & $F_n(T)$ over $\mathbb{F}_2$ & $F_n(T)$ over $\mathbb{F}_3$ & $F_n(T)$ over $\mathbb{F}_5$ \\
\midrule
1 & $1$ & $1$ & $1$ \\
2 & $T$ & $T$ & $T$ \\
3 & $(T + 1)^{2}$ & $(T^{2} + 1)$ & $(T + 2)(T + 3)$ \\
4 & $T^{3}$ & $T(T + 1)(T + 2)$ & $T(T^{2} + 2)$ \\
5 & $(T^{2} + T + 1)^{2}$ & $(T^{2} + T + 2)(T^{2} + 2 T + 2)$ & $(T + 1)^{2}(T + 4)^{2}$ \\
6 & $T(T + 1)^{4}$ & $T^{3}(T^{2} + 1)$ & $T(T + 2)(T + 3)(T^{2} + 3)$ \\
7 & $(T^{3} + T^{2} + 1)^{2}$ & $(T^{6} + 2 T^{4} + 1)$ & $(T^{3} + 2 T^{2} + 2 T + 2)(T^{3} + 3 T^{2} + 2 T + 3)$ \\
8 & $T^{7}$ & $T(T + 1)(T + 2)(T^{4} + T^{2} + 2)$ & $T(T^{2} + 2)(T^{4} + 4 T^{2} + 2)$ \\
9 & $(T + 1)^{2}(T^{3} + T + 1)^{2}$ & $(T^{2} + 1)^{4}$ & $(T + 2)(T + 3)(T^{3} + 3 T + 2)(T^{3} + 3 T + 3)$ \\
10 & $T(T^{2} + T + 1)^{4}$ & $T(T^{2} + T + 2)(T^{2} + 2 T + 2)(T^{4} + 2 T^{2} + 2)$ & $(T + 1)^{2}(T + 4)^{2}T^{5}$ \\
11 & $(T^{5} + T^{4} + T^{2} + T + 1)^{2}$ & $(T^{10} + T^{6} + 2 T^{4} + 1)$ & $(T^{5} + 2 T^{4} + 4 T^{3} + T^{2} + 3 T + 2)(T^{5} + 3 T^{4} + 4 T^{3} + 4 T^{2} + 3 T + 3)$ \\
12 & $T^{3}(T + 1)^{8}$ & $T^{3}(T + 1)^{3}(T + 2)^{3}(T^{2} + 1)$ & $T(T + 2)(T + 3)(T^{2} + 2)(T^{2} + 3)(T^{2} + 2 T + 4)(T^{2} + 3 T + 4)$ \\
13 & $(T^{6} + T^{5} + T^{4} + T + 1)^{2}$ & $(T^{6} + 2 T^{2} + 1)(T^{6} + 2 T^{4} + T^{2} + 1)$ & $(T^{2} + T + 1)(T^{2} + T + 2)(T^{2} + 2 T + 3)(T^{2} + 3 T + 3)(T^{2} + 4 T + 1)(T^{2} + 4 T + 2)$ \\
14 & $T(T^{3} + T^{2} + 1)^{4}$ & $T(T^{6} + T^{4} + 2 T^{2} + 1)(T^{6} + 2 T^{4} + 1)$ & $T(T^{3} + 2 T^{2} + 2 T + 2)(T^{3} + 3 T^{2} + 2 T + 3)(T^{6} + 2 T^{4} + 4 T^{2} + 2)$ \\
15 & $(T + 1)^{2}(T^{2} + T + 1)^{2}(T^{4} + T^{3} + 1)^{2}$ & $(T^{2} + 1)(T^{2} + T + 2)^{3}(T^{2} + 2 T + 2)^{3}$ & $(T + 1)^{2}(T + 4)^{2}(T + 2)^{5}(T + 3)^{5}$ \\
16 & $T^{15}$ & $T(T + 1)(T + 2)(T^{4} + T^{2} + 2)(T^{8} + 2 T^{6} + 2 T^{4} + T^{2} + 2)$ & $T(T^{2} + 2)(T^{4} + 4 T^{2} + 2)(T^{8} + 3 T^{6} + T^{2} + 2)$ \\
17 & $(T^{4} + T + 1)^{2}(T^{4} + T^{3} + T^{2} + T + 1)^{2}$ & $(T^{8} + T^{7} + 2 T^{6} + T^{5} + T^{4} + 2 T^{2} + T + 1)(T^{8} + 2 T^{7} + 2 T^{6} + 2 T^{5} + T^{4} + 2 T^{2} + 2 T + 1)$ & $(T^{8} + 2 T^{7} + 2 T^{6} + 2 T^{5} + 3 T + 1)(T^{8} + 3 T^{7} + 2 T^{6} + 3 T^{5} + 2 T + 1)$ \\
18 & $T(T + 1)^{4}(T^{3} + T + 1)^{4}$ & $T^{9}(T^{2} + 1)^{4}$ & $T(T + 2)(T + 3)(T^{2} + 3)(T^{3} + 3 T + 2)(T^{3} + 3 T + 3)(T^{6} + T^{4} + 4 T^{2} + 3)$ \\
19 & $(T^{9} + T^{8} + T^{6} + T^{5} + T^{4} + T + 1)^{2}$ & $(T^{18} + 2 T^{16} + 2 T^{12} + 2 T^{10} + 1)$ & $(T^{9} + 2 T^{8} + 3 T^{7} + 4 T^{6} + T^{5} + 2)(T^{9} + 3 T^{8} + 3 T^{7} + T^{6} + T^{5} + 3)$ \\
20 & $T^{3}(T^{2} + T + 1)^{8}$ & $T(T + 1)(T + 2)(T^{2} + T + 2)(T^{2} + 2 T + 2)(T^{4} + 2 T^{2} + 2)(T^{4} + T^{3} + 2)(T^{4} + 2 T^{3} + 2)$ & $(T + 1)^{2}(T + 4)^{2}T^{5}(T^{2} + 2)^{5}$ \\
\bottomrule
\end{tabular}
\caption{Factorization of $F_n(T)$ over $\mathbb{F}_2$, $\mathbb{F}_3$, and $\bF_5$}
\label{tab:fibo}
\end{table}



\begin{table}
\centering
\begin{tabular}{c|p{0.45\textwidth}|p{0.45\textwidth}}
\toprule
$n$ & $U_n(T)$ over $\mathbb{F}_3$ & $U_n(T)$ over $\mathbb{F}_5$ \\
\midrule
1 & $2T$ & $2T$ \\
2 & $(T + 1)(T + 2)$ & $4(T + 2)(T + 3)$ \\
3 & $2T(T^{2} + 1)$ & $3T(T^{2} + 2)$ \\
4 & $(T^{2} + T + 2)(T^{2} + 2 T + 2)$ & $(T + 1)^{2}(T + 4)^{2}$ \\
5 & $2(T + 1)(T + 2)T^{3}$ & $2T(T + 2)(T + 3)(T^{2} + 3)$ \\
6 & $(T^{3} + T^{2} + T + 2)(T^{3} + 2 T^{2} + T + 1)$ & $4(T^{3} + 2 T^{2} + 2 T + 2)(T^{3} + 3 T^{2} + 2 T + 3)$ \\
7 & $2T(T^{2} + 1)(T^{4} + 2 T^{2} + 2)$ & $3T(T^{2} + 2)(T^{4} + 4 T^{2} + 2)$ \\
8 & $(T + 1)^{4}(T + 2)^{4}$ & $(T + 2)(T + 3)(T^{3} + 3 T + 2)(T^{3} + 3 T + 3)$ \\
9 & $2T(T^{2} + T + 2)(T^{2} + 2 T + 2)(T^{4} + T^{2} + 2)$ & $2(T + 1)^{2}(T + 4)^{2}T^{5}$ \\
10 & $(T^{5} + T^{4} + 2 T^{3} + 1)(T^{5} + 2 T^{4} + 2 T^{3} + 2)$ & $4(T^{5} + 2 T^{4} + 4 T^{3} + T^{2} + 3 T + 2)(T^{5} + 3 T^{4} + 4 T^{3} + 4 T^{2} + 3 T + 3)$ \\
11 & $2(T + 1)(T + 2)T^{3}(T^{2} + 1)^{3}$ & $3T(T + 2)(T + 3)(T^{2} + 2)(T^{2} + 3)(T^{2} + 2 T + 4)(T^{2} + 3 T + 4)$ \\
12 & $(T^{3} + 2 T + 1)(T^{3} + 2 T + 2)(T^{3} + T^{2} + 2 T + 1)(T^{3} + 2 T^{2} + 2 T + 2)$ & $(T^{2} + T + 1)(T^{2} + T + 2)(T^{2} + 2 T + 3)(T^{2} + 3 T + 3)(T^{2} + 4 T + 1)(T^{2} + 4 T + 2)$ \\
13 & $2T(T^{3} + T^{2} + 2)(T^{3} + T^{2} + T + 2)(T^{3} + 2 T^{2} + 1)(T^{3} + 2 T^{2} + T + 1)$ & $2T(T^{3} + 2 T^{2} + 2 T + 2)(T^{3} + 3 T^{2} + 2 T + 3)(T^{6} + 2 T^{4} + 4 T^{2} + 2)$ \\
14 & $(T + 1)(T + 2)(T^{2} + T + 2)^{3}(T^{2} + 2 T + 2)^{3}$ & $4(T + 1)^{2}(T + 4)^{2}(T + 2)^{5}(T + 3)^{5}$ \\
15 & $2T(T^{2} + 1)(T^{4} + 2 T^{2} + 2)(T^{8} + T^{6} + 2 T^{4} + 2 T^{2} + 2)$ & $3T(T^{2} + 2)(T^{4} + 4 T^{2} + 2)(T^{8} + 3 T^{6} + T^{2} + 2)$ \\
16 & $(T^{8} + T^{7} + 2 T^{6} + T^{3} + 2 T^{2} + 2 T + 1)(T^{8} + 2 T^{7} + 2 T^{6} + 2 T^{3} + 2 T^{2} + T + 1)$ & $(T^{8} + 2 T^{7} + 2 T^{6} + 2 T^{5} + 3 T + 1)(T^{8} + 3 T^{7} + 2 T^{6} + 3 T^{5} + 2 T + 1)$ \\
17 & $2(T + 1)^{4}(T + 2)^{4}T^{9}$ & $2T(T + 2)(T + 3)(T^{2} + 3)(T^{3} + 3 T + 2)(T^{3} + 3 T + 3)(T^{6} + T^{4} + 4 T^{2} + 3)$ \\
18 & $(T^{9} + T^{8} + T^{7} + 2 T^{6} + T^{3} + 2 T^{2} + 2 T + 1)(T^{9} + 2 T^{8} + T^{7} + T^{6} + T^{3} + T^{2} + 2 T + 2)$ & $4(T^{9} + 2 T^{8} + 3 T^{7} + 4 T^{6} + T^{5} + 2)(T^{9} + 3 T^{8} + 3 T^{7} + T^{6} + T^{5} + 3)$ \\
19 & $2T(T^{2} + 1)(T^{2} + T + 2)(T^{2} + 2 T + 2)(T^{4} + T^{2} + 2)(T^{4} + T^{3} + T^{2} + 2 T + 2)(T^{4} + 2 T^{3} + T^{2} + T + 2)$ & $3(T + 1)^{2}(T + 4)^{2}T^{5}(T^{2} + 2)^{5}$ \\
20 & $(T + 1)(T + 2)(T^{3} + T^{2} + T + 2)^{3}(T^{3} + 2 T^{2} + T + 1)^{3}$ & $(T + 2)(T + 3)(T^{3} + T + 1)(T^{3} + T + 4)(T^{3} + 2 T^{2} + 1)(T^{3} + 2 T^{2} + 2 T + 2)(T^{3} + 3 T^{2} + 4)(T^{3} + 3 T^{2} + 2 T + 3)$ \\
\bottomrule
\end{tabular}
\caption{Factorization of $U_n(T)$ over $\mathbb{F}_3$ and $\mathbb{F}_5$}
\label{tab:chebyshev2}
\end{table}



\begin{table}
\centering
\begin{tabular}{c|p{0.30\textwidth}|p{0.30\textwidth}|p{0.30\textwidth}}
\toprule
$n$ & $L_n(T)$ over $\mathbb{F}_2$ & $L_n(T)$ over $\mathbb{F}_3$ & $L_n(T)$ over $\mathbb{F}_5$ \\
\midrule
1 & $T$ & $T$ & $T$ \\
2 & $T^{2}$ & $(T + 1)(T + 2)$ & $(T^{2} + 2)$ \\
3 & $T(T + 1)^{2}$ & $T^{3}$ & $T(T^{2} + 3)$ \\
4 & $T^{4}$ & $(T^{4} + T^{2} + 2)$ & $(T^{4} + 4 T^{2} + 2)$ \\
5 & $T(T^{2} + T + 1)^{2}$ & $T(T^{4} + 2 T^{2} + 2)$ & $T^{5}$ \\
6 & $T^{2}(T + 1)^{4}$ & $(T + 1)^{3}(T + 2)^{3}$ & $(T^{2} + 2)(T^{2} + 2 T + 4)(T^{2} + 3 T + 4)$ \\
7 & $T(T^{3} + T^{2} + 1)^{2}$ & $T(T^{6} + T^{4} + 2 T^{2} + 1)$ & $T(T^{6} + 2 T^{4} + 4 T^{2} + 2)$ \\
8 & $T^{8}$ & $(T^{8} + 2 T^{6} + 2 T^{4} + T^{2} + 2)$ & $(T^{8} + 3 T^{6} + T^{2} + 2)$ \\
9 & $T(T + 1)^{2}(T^{3} + T + 1)^{2}$ & $T^{9}$ & $T(T^{2} + 3)(T^{6} + T^{4} + 4 T^{2} + 3)$ \\
10 & $T^{2}(T^{2} + T + 1)^{4}$ & $(T + 1)(T + 2)(T^{4} + T^{3} + 2)(T^{4} + 2 T^{3} + 2)$ & $(T^{2} + 2)^{5}$ \\
11 & $T(T^{5} + T^{4} + T^{2} + T + 1)^{2}$ & $T(T^{5} + T^{4} + T + 2)(T^{5} + 2 T^{4} + T + 1)$ & $T(T^{5} + 2 T^{4} + T^{2} + 4 T + 3)(T^{5} + 3 T^{4} + 4 T^{2} + 4 T + 2)$ \\
12 & $T^{4}(T + 1)^{8}$ & $(T^{4} + T^{2} + 2)^{3}$ & $(T^{4} + 2)(T^{4} + 3 T^{2} + 3)(T^{4} + 4 T^{2} + 2)$ \\
13 & $T(T^{6} + T^{5} + T^{4} + T + 1)^{2}$ & $T(T^{3} + T^{2} + T + 2)(T^{3} + T^{2} + 2 T + 1)(T^{3} + 2 T^{2} + T + 1)(T^{3} + 2 T^{2} + 2 T + 2)$ & $T(T^{4} + 3)(T^{4} + T^{2} + 2)(T^{4} + 2 T^{2} + 3)$ \\
14 & $T^{2}(T^{3} + T^{2} + 1)^{4}$ & $(T + 1)(T + 2)(T^{3} + 2 T + 1)(T^{3} + 2 T + 2)(T^{3} + T^{2} + 2)(T^{3} + 2 T^{2} + 1)$ & $(T^{2} + 2)(T^{6} + T^{5} + 4 T^{4} + 3 T^{3} + 4 T^{2} + 3 T + 1)(T^{6} + 4 T^{5} + 4 T^{4} + 2 T^{3} + 4 T^{2} + 2 T + 1)$ \\
15 & $T(T + 1)^{2}(T^{2} + T + 1)^{2}(T^{4} + T^{3} + 1)^{2}$ & $T^{3}(T^{4} + 2 T^{2} + 2)^{3}$ & $T^{5}(T^{2} + 3)^{5}$ \\
16 & $T^{16}$ & $(T^{16} + T^{14} + 2 T^{12} + T^{10} + T^{2} + 2)$ & $(T^{16} + T^{14} + 4 T^{12} + 2 T^{10} + 2 T^{6} + T^{4} + 4 T^{2} + 2)$ \\
17 & $T(T^{4} + T + 1)^{2}(T^{4} + T^{3} + T^{2} + T + 1)^{2}$ & $T(T^{16} + 2 T^{14} + 2 T^{12} + T^{10} + 2 T^{8} + 2)$ & $T(T^{16} + 2 T^{14} + 4 T^{12} + 2 T^{10} + 2 T^{6} + 4 T^{4} + 4 T^{2} + 2)$ \\
18 & $T^{2}(T + 1)^{4}(T^{3} + T + 1)^{4}$ & $(T + 1)^{9}(T + 2)^{9}$ & $(T^{2} + 2)(T^{2} + 2 T + 4)(T^{2} + 3 T + 4)(T^{6} + T^{4} + 2 T^{3} + 4 T^{2} + T + 4)(T^{6} + T^{4} + 3 T^{3} + 4 T^{2} + 4 T + 4)$ \\
19 & $T(T^{9} + T^{8} + T^{6} + T^{5} + T^{4} + T + 1)^{2}$ & $T(T^{18} + T^{16} + 2 T^{14} + 2 T^{12} + T^{10} + 2 T^{8} + 1)$ & $T(T^{9} + T^{8} + T^{5} + 2 T^{4} + 2 T^{3} + 2 T^{2} + T + 4)(T^{9} + 4 T^{8} + T^{5} + 3 T^{4} + 2 T^{3} + 3 T^{2} + T + 1)$ \\
20 & $T^{4}(T^{2} + T + 1)^{8}$ & $(T^{4} + T + 2)(T^{4} + 2 T + 2)(T^{4} + T^{2} + 2)(T^{4} + T^{3} + T^{2} + T + 1)(T^{4} + 2 T^{3} + T^{2} + 2 T + 1)$ & $(T^{4} + 4 T^{2} + 2)^{5}$ \\
\bottomrule
\end{tabular}
\caption{Factorization of $L_n(T)$ over $\mathbb{F}_2$,  $\mathbb{F}_3$, and $\mathbb{F}_5$}
\label{tab:lucas}
\end{table}

\begin{table}
\centering
\begin{tabular}{c|p{0.45\textwidth}|p{0.45\textwidth}}
\toprule
$n$ & $T_n(T)$ over $\mathbb{F}_3$ & $T_n(T)$ over $\mathbb{F}_5$ \\
\midrule
1 & $T$ & $T$ \\
2 & $2(T^{2} + 1)$ & $2(T^{2} + 2)$ \\
3 & $T^{3}$ & $4T(T^{2} + 3)$ \\
4 & $2(T^{4} + 2 T^{2} + 2)$ & $3(T^{4} + 4 T^{2} + 2)$ \\
5 & $T(T^{4} + T^{2} + 2)$ & $T^{5}$ \\
6 & $2(T^{2} + 1)^{3}$ & $2(T^{2} + 2)(T^{2} + 2 T + 4)(T^{2} + 3 T + 4)$ \\
7 & $T(T^{3} + T^{2} + 2)(T^{3} + 2 T^{2} + 1)$ & $4T(T^{6} + 2 T^{4} + 4 T^{2} + 2)$ \\
8 & $2(T^{8} + T^{6} + 2 T^{4} + 2 T^{2} + 2)$ & $3(T^{8} + 3 T^{6} + T^{2} + 2)$ \\
9 & $T^{9}$ & $T(T^{2} + 3)(T^{6} + T^{4} + 4 T^{2} + 3)$ \\
10 & $2(T^{2} + 1)(T^{4} + T^{3} + T^{2} + 2 T + 2)(T^{4} + 2 T^{3} + T^{2} + T + 2)$ & $2(T^{2} + 2)^{5}$ \\
11 & $T(T^{10} + T^{8} + 2 T^{6} + T^{4} + T^{2} + 1)$ & $4T(T^{5} + 2 T^{4} + T^{2} + 4 T + 3)(T^{5} + 3 T^{4} + 4 T^{2} + 4 T + 2)$ \\
12 & $2(T^{4} + 2 T^{2} + 2)^{3}$ & $3(T^{4} + 2)(T^{4} + 3 T^{2} + 3)(T^{4} + 4 T^{2} + 2)$ \\
13 & $T(T^{6} + 2 T^{2} + 1)(T^{6} + 2 T^{4} + 1)$ & $T(T^{4} + 3)(T^{4} + T^{2} + 2)(T^{4} + 2 T^{2} + 3)$ \\
14 & $2(T^{2} + 1)(T^{6} + T^{4} + 2 T^{2} + 1)(T^{6} + 2 T^{4} + T^{2} + 1)$ & $2(T^{2} + 2)(T^{6} + T^{5} + 4 T^{4} + 3 T^{3} + 4 T^{2} + 3 T + 1)(T^{6} + 4 T^{5} + 4 T^{4} + 2 T^{3} + 4 T^{2} + 2 T + 1)$ \\
15 & $T^{3}(T^{4} + T^{2} + 2)^{3}$ & $4T^{5}(T^{2} + 3)^{5}$ \\
16 & $2(T^{16} + 2 T^{14} + 2 T^{12} + 2 T^{10} + 2 T^{2} + 2)$ & $3(T^{16} + T^{14} + 4 T^{12} + 2 T^{10} + 2 T^{6} + T^{4} + 4 T^{2} + 2)$ \\
17 & $T(T^{16} + T^{14} + 2 T^{12} + 2 T^{10} + 2 T^{8} + 2)$ & $T(T^{16} + 2 T^{14} + 4 T^{12} + 2 T^{10} + 2 T^{6} + 4 T^{4} + 4 T^{2} + 2)$ \\
18 & $2(T^{2} + 1)^{9}$ & $2(T^{2} + 2)(T^{2} + 2 T + 4)(T^{2} + 3 T + 4)(T^{6} + T^{4} + 2 T^{3} + 4 T^{2} + T + 4)(T^{6} + T^{4} + 3 T^{3} + 4 T^{2} + 4 T + 4)$ \\
19 & $T(T^{9} + T^{8} + T^{6} + 2 T^{5} + T^{4} + 2 T^{3} + T^{2} + 2 T + 2)(T^{9} + 2 T^{8} + 2 T^{6} + 2 T^{5} + 2 T^{4} + 2 T^{3} + 2 T^{2} + 2 T + 1)$ & $4T(T^{9} + T^{8} + T^{5} + 2 T^{4} + 2 T^{3} + 2 T^{2} + T + 4)(T^{9} + 4 T^{8} + T^{5} + 3 T^{4} + 2 T^{3} + 3 T^{2} + T + 1)$ \\
20 & $2(T^{4} + 2 T^{2} + 2)(T^{4} + T^{3} + 2 T + 1)(T^{4} + T^{3} + 2 T^{2} + 2 T + 2)(T^{4} + 2 T^{3} + T + 1)(T^{4} + 2 T^{3} + 2 T^{2} + T + 2)$ & $3(T^{4} + 4 T^{2} + 2)^{5}$ \\
\bottomrule
\end{tabular}
\caption{Factorization of $T_n(T)$ over $\mathbb{F}_3$ and $\mathbb{F}_5$}
\label{tab:chebyshev1}
\end{table}


\begin{table}
\centering
\begin{tabular}{c|p{0.45\textwidth}|p{0.45\textwidth}}
\toprule
$n$ & $J_n(T)$ over $\mathbb{F}_3$ & $J_n(T)$ over $\mathbb{F}_5$ \\
\midrule
1 & $1$ & $1$ \\
2 & $1$ & $1$ \\
3 & $2(T + 2)$ & $2(T + 3)$ \\
4 & $(T + 1)$ & $4(T + 4)$ \\
5 & $(T^{2} + 1)$ & $4(T + 2)^{2}$ \\
6 & $2(T + 2)$ & $2(T + 1)(T + 3)$ \\
7 & $2(T^{3} + 2 T + 2)$ & $3(T^{3} + 3 T^{2} + 2)$ \\
8 & $2(T + 1)(T^{2} + T + 2)$ & $2(T + 4)(T^{2} + T + 2)$ \\
9 & $(T + 2)^{4}$ & $(T + 3)(T^{3} + 2 T^{2} + 4 T + 2)$ \\
10 & $2(T^{2} + 1)(T^{2} + 2 T + 2)$ & $4(T + 2)^{2}$ \\
11 & $2(T^{5} + 2 T^{3} + 2 T^{2} + 2)$ & $2(T^{5} + T^{2} + 4 T + 3)$ \\
12 & $2(T + 2)(T + 1)^{3}$ & $2(T + 1)(T + 3)(T + 4)(T^{2} + 2 T + 4)$ \\
13 & $(T^{3} + T^{2} + 2)(T^{3} + 2 T^{2} + 2 T + 2)$ & $4(T^{2} + T + 1)(T^{2} + 3 T + 4)(T^{2} + 4 T + 1)$ \\
14 & $(T^{3} + 2 T + 2)(T^{3} + T^{2} + T + 2)$ & $3(T^{3} + T^{2} + 4 T + 1)(T^{3} + 3 T^{2} + 2)$ \\
15 & $2(T + 2)(T^{2} + 1)^{3}$ & $3(T + 2)^{2}(T + 3)^{5}$ \\
16 & $(T + 1)(T^{2} + T + 2)(T^{4} + T^{3} + T^{2} + 2 T + 2)$ & $4(T + 4)(T^{2} + T + 2)(T^{4} + 4 T^{3} + 3 T + 3)$ \\
17 & $(T^{8} + 2 T^{3} + T^{2} + 1)$ & $(T^{8} + 3 T^{7} + 4 T^{5} + 3 T^{3} + 4 T^{2} + 1)$ \\
18 & $(T + 2)^{4}$ & $4(T + 1)(T + 3)(T^{3} + 2 T^{2} + 4 T + 2)(T^{3} + 4 T^{2} + 3 T + 4)$ \\
19 & $2(T^{9} + T^{4} + 2 T^{3} + 2 T + 2)$ & $2(T^{9} + 3 T^{6} + 2 T^{5} + 3 T^{4} + 2 T + 3)$ \\
20 & $2(T + 1)(T^{2} + 1)(T^{2} + 2 T + 2)(T^{4} + T^{2} + T + 1)$ & $(T + 2)^{2}(T + 4)^{5}$ \\
21 & $(T + 2)(T^{3} + 2 T + 2)^{3}$ & $4(T + 3)(T^{3} + 2 T^{2} + T + 3)(T^{3} + 2 T^{2} + 2 T + 3)(T^{3} + 3 T^{2} + 2)$ \\
22 & $2(T^{5} + 2 T^{3} + 2 T^{2} + 2)(T^{5} + T^{4} + T^{3} + 2 T^{2} + T + 1)$ & $4(T^{5} + T^{2} + 4 T + 3)(T^{5} + 3 T^{3} + 3 T^{2} + T + 3)$ \\
23 & $2(T^{11} + T^{9} + 2 T^{6} + 2 T^{5} + 2 T^{2} + 2)$ & $3(T^{11} + 3 T^{10} + T^{8} + 4 T^{6} + 2 T^{5} + 4 T^{3} + 4 T + 2)$ \\
24 & $(T + 2)(T + 1)^{3}(T^{2} + T + 2)^{3}$ & $(T + 1)(T + 3)(T + 4)(T^{2} + 2)(T^{2} + T + 2)(T^{2} + 2 T + 4)(T^{2} + 3 T + 3)$ \\
25 & $(T^{2} + 1)(T^{10} + T^{8} + 2 T^{7} + 2 T^{6} + 2 T^{5} + T^{4} + T^{3} + 2 T^{2} + T + 1)$ & $(T + 2)^{12}$ \\
26 & $(T^{3} + 2 T + 1)(T^{3} + T^{2} + 2)(T^{3} + 2 T^{2} + 1)(T^{3} + 2 T^{2} + 2 T + 2)$ & $3(T^{2} + 3)(T^{2} + T + 1)(T^{2} + 2 T + 3)(T^{2} + 3 T + 4)(T^{2} + 4 T + 1)(T^{2} + 4 T + 2)$ \\
27 & $2(T + 2)^{13}$ & $2(T + 3)(T^{3} + 2 T^{2} + 4 T + 2)(T^{9} + 3 T^{8} + 2 T^{6} + 2 T^{5} + T^{4} + 4 T^{3} + 3 T + 3)$ \\
\bottomrule
\end{tabular}
\caption{Factorization of $J_n(T)$ over $\bF_3$ and $\mathbb{F}_5$}
\label{tab:jacobsthal}
\end{table}